\scanpage{217}
\mspage{72}

\[
 \left(\bigcup A_i\right)^{\unsure{\sharp}}=\left(\bigcup_i A_i^\sharp\right)^\sharp\ ;
\]

Pour tout $A\inserted{\subset P}$, $A^\sharp$ est la plus petite partie $R$-close de $P$ contenant $A$, et de $=$ pour tt $B\subset Q$, $B^\flat$ est la plus petite partie $R$-close de $Q$ contenant $B$. On a de plus (c'était $=$ à dire avant de dire que $A\mapsto A^\sharp$ est croissante)

\[
\left\{\begin{aligned}
 A\subset A'&\Longrightarrow A'^{\unsure{\cup|\smile}}\subset A^\cup\\
 B\subset B'&\Longrightarrow B'^{\unsure{\cap|\frown}}\subset B^\cap,
\end{aligned}\right.
\tag*{(5.11)}
\]

i.e. les applications

\diagram{\begin{tikzcd}[ampersand replacement=\&,column sep=large]
\text{(5.12)} \& \unsure{\mathfrak P|\mathcal P}(\unsure{p|P})\arrow[r,shift left,"\unsure{B|\beta}"]
 \& \mathfrak P(Q)\arrow[l,shift left,"\alpha"]
 \& \begin{aligned}A&\unsure{\vdash|\mapsto} A^\cup\\ B&\vdash B^\cap\end{aligned}
\end{tikzcd}}

sont décroissantes \unsure{(ce|Ce} qui implique que les deux composés $\alpha\beta=(A\mapsto A^\sharp)$ et $\beta\alpha=(B\mapsto B^\flat)$ sont croissantes). Leurs restrictions aux sous-ensembles $\mathfrak P_R(P)$ et $\mathfrak P_{\unsure{{}^{s}R|{}^{t}R}}(Q)$ des parties $R$-closes de $P$ resp. de $Q$, sont \deleted{\unsure{donc} telles}, par définition de celles-ci comme

\[
\begin{aligned}
 \mathfrak P_R(P)&=\{A\subset P\mid A^\sharp(\overset{\mathrm{def}}{=}\alpha\beta(A))=A\}\\
 \mathfrak P_R(Q)&=\{B\subset Q\mid B^\flat(\overset{\mathrm{def}}{=}\beta\alpha(\unsure{B,|B}))=B\},
\end{aligned}
\]

inverses l'une de l'autre. On trouve ainsi

\uline{Proposition} 5.1\quad Soient $P,Q$ deux ensembles, $R\inserted{\subset P\times Q}$ une relation. Alors les applications $\alpha,\beta$ de (5.12.) \unsure{induisent} des bijections, inverses l' 
% archive sheet number: 217
